\begin{textsample}{POS Dim 4 – 2012\_c6 – Score 3.00 – 2012\_c6/enade\_2012\_ci\_econ\_1bf88b0c63.txt}  \label{ex:f4_pos_121}
CAPÍTULO 3 ANÁLISE TÉCNICA DA PROVA Este capítulo tem por objetivo apresentar o desempenho dos estudantes concluintes de Ciências Econômicas no ENADE/2012.
Para isso, foram calculadas as estatísticas básicas da prova em seu todo, bem como as estatísticas dos componentes relacionadas à Formação Geral, ao de Conhecimento Específico da Área e das questões discursivas isoladamente.
Nas tabelas, são apresentados o tamanho da população inscrita e de presentes, e as seguintes estatísticas das notas9 : média do desempenho na prova, erro padrão da média, desvio padrão, nota mínima, mediana e nota máxima.
Tais estatísticas contemplam o total de estudantes concluintes da Área de Ciências Econômicas presentes à prova do ENADE 2012, tendo -se em vista as seguintes agregações : ( a ) as Grandes Regiões e o país como um todo ; ( b ) a Categoria Administrativa ; e ( c ) a Organização Acadêmica.
Em relação aos gráficos de distribuição de notas, o intervalo considerado foi de 10 unidades, aberto à esquerda e fechado à direita, com exceção do primeiro intervalo, [ 0 ; 10 ], fechado em ambos os extremos.
Para os gráficos de distribuição das notas das questões discursivas, foram consideradas mais duas categorias : questão em branco e nota zero. 3.1 ESTATÍSTICAS BÁSICAS DA PROVA 3.1.1 Estatísticas Básicas Gerais A Tabela 3.1 apresenta as estatísticas básicas da prova por grande Região.
A população total de inscritos foi de 10.528.
Destes, 8.257 estiveram presentes, sendo 21,6 \% o índice de não comparecimento.
A região de maior abstenção foi a Norte ( 38,2 \% ) e a de menor abstenção foi a Sul ( 17,8 \% ).
A média das notas da prova como um todo ( nas seções seguintes serão analisados os componentes de Formação Geral e de Conhecimento Específico ) foi 31,1, sendo que os alunos da região Nordeste obtiveram a média mais baixa ( 29,3 ) e os da região Centro-Oeste obtiveram a média mais alta ( 31,8 ).
As demais médias foram : 31,7 na região Sudeste, 31,5 na região Sul e 30,4 na região Norte.
O desvio padrão para o Brasil como um todo foi 12,5, 9 Essas estatísticas e outras estão definidas no Capítulo 1. sendo o maior desvio padrão encontrado na região Sudeste ( 13,6 ) e o menor na região Norte ( 9,7 ), indicando uma menor dispersão das notas desta última região.
A região que obteve a maior nota máxima foi a Sudeste ( 86,9 ), ao passo que a região que atingiu a menor nota máxima foi a Nordeste ( 69,2 ).
A mediana do Brasil como um todo foi 30,6, sendo a maior mediana obtida na região Sul ( 31,1 ) e a menor obtida na Nordeste ( 29,5 ).
A nota mínima foi zero em todas as regiões do Brasil, sem exceção.
O comportamento das notas dos estudantes de todo o Brasil pode ser observado no Gráfico 3.1 que apresenta um histograma com a distribuição das mesmas.
Essa é uma distribuição unimodal com moda no intervalo ( 20 ; 30 ].
Destaca -se, também, o intervalo ( 30 ; 40 ] com frequência relativa muito próxima à modal.
O coeficiente de assimetria da distribuição das notas é positivo ( 0,18 ), indicando uma assimetria à direita.
As distribuições por Grande Região também apresentam assimetria positiva, concentração pouco maior do lado esquerdo do histograma e mais espalhada do lado direito.
A única exceção é a região Nordeste, com coeficiente de assimetria negativo ( -0,21 ), a distribuição tem comportamento oposto : concentração à direita e cauda um pouco mais longa à esquerda.
Os Gráficos 3.2, 3.3 e 3.4 apresentam informações referentes à média da nota final dos participantes, desagregando os resultados de acordo com, respectivamente, as Grandes Regiões do país, a Categoria Administrativa e a Organização Acadêmica.
Os gráficos apresentam o valor da média das notas como uma barra e os extremos do intervalo de confiança de 95 \% como linhas verticais unidas por uma linha horizontal na forma da letra H maiúscula.
Considerando -se o gráfico de notas segundo Grande Região ( Gráfico 3.2 ), observa -se que existe diferença estatisticamente significativa ao nível de 95 \% entre a maior média, obtida na região Centro-Oeste ( 31,8 ) e a menor, obtida na região Nordeste ( 29,3 ).
Levando -se em conta os estudantes por Categorias Administrativas ( Gráfico 3.3 ), observa -se que não há diferença entre as médias das notas das IES Públicas e Privadas ( 31,1 ).
Tendo como foco o Gráfico 3.4, que apresenta as notas médias das provas segundo Organização Acadêmica, constata -se que existe diferença estatisticamente significativa ao nível de 95 \% nas médias das notas dos estudantes provenientes de Faculdades em relação às médias de Centros Universitários e Universidades.
A maior média foi obtida pelos estudantes de Faculdades ( 33,4 ), e a menor, pelos de Centros Universitários ( 30,2 ). 3.1.2 Estatísticas Básicas no Componente de Formação Geral A Tabela 3.2 apresenta as estatísticas básicas em relação ao componente da prova que avalia a Formação Geral dos estudantes concluintes.
Os alunos de todo Brasil obtiveram desempenho médio de 43,3.
Quanto à variabilidade, o desvio padrão das notas dos estudantes do Brasil como um todo foi 18,3.
A maior média foi obtida na região Centro-Oeste ( 46,1 ), e a menor, na região Nordeste ( 42,6 ).
As demais médias foram : 44,8 na região Norte, 44,4 na região Sul e 42,7 na região Sudeste.
Já o maior desvio padrão foi obtido na região Sudeste ( 19,4 ) e o menor na região Sul ( 16,4 ).
Os demais desvios padrões foram : 17,9 na região Centro-Oeste e 17,8 nas regiões Norte e Nordeste.
A maior nota no componente de Formação Geral da prova do ENADE foi obtida por pelo menos um aluno da região Sudeste ( 92,0 ) enquanto que a menor nota máxima foi obtida na região Norte ( 81,5 ).
Nas outras regiões as notas máximas foram : 90,0 na região Nordeste, 87,0 na região Centro-Oeste e 84,5 na região Sul.
A mediana do Brasil como um todo foi 45,0, sendo esta igual a menor mediana, encontrada nas regiões Nordeste, Sudeste e Sul.
A maior mediana foi encontrada na região Centro-Oeste ( 47,0 ).
A nota mínima de Formação Geral foi zero em todas as regiões, sem exceção.
O Gráfico 3.5 propicia a avaliação do desempenho dos estudantes no componente de Formação Geral a partir do histograma da distribuição das notas correspondentes.
A distribuição é unimodal, com moda em ( 40 ; 50 ], enquanto na prova como um todo a moda foi alcançada no intervalo ( 20 ; 30 ].
Nota -se, ainda, que no gráfico 3.5 as notas apresentam uma maior dispersão do que no Gráfico 3.1 ( distribuição das notas da prova ), confirmado pela comparação dos desvios padrões : 12,5 para a nota da prova como um todo e 18,3 para o componente de Formação Geral.
Para o componente de Formação Geral, o coeficiente de assimetria da distribuição das notas dos estudantes, diferentemente da prova como um todo, é negativo ( -0,42 ), indicando uma assimetria à esquerda.
Em todas as Grandes Regiões os histogramas também possuem assimetria levemente negativa ( entre -0,53 para a região Centro-Oeste e -0,26 para a região Norte ).
Nos Gráficos 3.6, 3.7 e 3.8 são apresentadas as informações referentes ao desempenho dos concluintes no componente de Formação Geral, em diferentes agregações : Grande Região do país, Categoria Administrativa e Organização Acadêmica.
Observa -se pelo Gráfico 3.6 que existe diferença estatisticamente significativa entre a maior e as duas menores médias das notas no Componente de Formação Geral, a maior obtida na região Centro-Oeste ( 46,1 ), e as duas menores nas regiões Nordeste ( 42,6 ) e Sudeste ( 42,7 ).
Vemos que o intervalo de confiança mais largo é o da região Norte ; já o intervalo mais estreito é observado na região Sudeste.
Este fato está relacionado, também, com o tamanho da população envolvida, menor na região Norte do que na Sudeste.
No Gráfico 3.7, que representa as notas médias no Componente de Formação Geral segundo Categoria Administrativa do país, observa -se que não existe diferença estatisticamente significativa entre as médias das IES públicas ( 43,0 ) e privadas ( 43,7 ), como ocorreu para as notas da prova como um todo, cujos valores eram idênticos para as duas categorias.
Considerando -se o tipo de Organização Acadêmica, nota -se, no Gráfico 3.8, uma diferença estatisticamente significativa entre a maior e a menor média.
Nas Faculdades ( 46,7 ), a média é mais elevada do que em Centros Universitários ( 45,1 ) e Universidades ( 42,6 ). 3.1.3 Estatísticas Básicas do Componente de Conhecimento Específico A Tabela 3.3 apresenta as estatísticas básicas referentes ao componente de Conhecimento Específico da Área de Ciências Econômicas.
A média do desempenho dos alunos do Brasil como um todo foi 27,0.
A maior média foi obtida na região Sudeste ( 28,0 ), e a menor, na região Nordeste ( 24,8 ).
As demais médias foram : 27,1 na região Sul, 27,0 na região Centro-Oeste e 25,7 na região Norte.
Quanto à variabilidade das notas, o desvio padrão do Brasil como um todo foi 13,0, sendo o maior desvio padrão observado na região Centro-Oeste ( 14,3 ) e o menor na região Norte ( 10,2 ).
Os demais desvios foram : 14,0 da região Sudeste, 11,8 na região Sul e 11,1 da região Nordeste.
A mediana das notas dos estudantes de todo o Brasil foi 25,3.
A maior mediana ocorreu nas regiões Sudeste e Sul ( 26,1 ) e a menor nas regiões Norte, Nordeste e Centro-Oeste ( 24,8 ).
A nota máxima do Brasil como um todo foi 92,9, sendo obtida por pelo menos um aluno da região Sudeste.
As demais notas máximas foram : 88,9 na região Centro-Oeste, 77,3 na região Sul, 72,2 na região Norte e 70,2 na região Nordestee.
A nota mínima foi zero em todas as regiões, sem exceção.
Assim como os Gráficos 3.1 e 3.5, o Gráfico 3.9, apresentado a seguir, proporciona uma avaliação do desempenho de concluintes em relação ao componente de Conhecimento Específico com um histograma da distribuição das notas correspondentes.
Dentre as três distribuições apresentadas, esta é a mais concentrada nas notas baixas.
Esta também é uma distribuição unimodal e o grupo modal é o ( 20 ; 30 ], abaixo do grupo modal para a Formação Geral.
O coeficiente de assimetria da distribuição das notas do componente de Conhecimento Específico é positivo ( 0,60 ), indicando assimetria à direita.
Em todas as Grandes Regiões os histogramas também possuem assimetria positiva ( entre 0,27 para a região Nordeste e 1,12 para a região Centro-Oeste ).
Os Gráficos 3.10, 3.11 e 3.12 apresentam, respectivamente, uma comparação dos resultados em relação à Grande Região do país, à Categoria Administrativa e à Organização Acadêmica, agora levando em conta o desempenho de alunos no componente de Conhecimento Específico da prova.
Pelo Gráfico 3.10, observa -se que existe diferença estatisticamente significativa entre a maior média das notas no componente de Conhecimento Específico, obtida na região Sudeste ( 28,0 ) e as duas menores, obtidas nas regiões Nordeste ( 24,8 ) e Norte ( 25,7 ).
Quanto à Categoria Administrativa ( Gráfico 3.11 ), observa -se um comportamento similar àquele da parte de Formação Geral, ou seja, não existe diferença estatisticamente significativa entre as médias das IES Públicas e Privadas.
Além disso, a maior média foi obtida por alunos de IES Públicas de ensino ( 27,1 ), ao contrário da parte de Formação Geral.
Quanto ao Gráfico 3.12, observa -se, mais uma vez, que existe diferença estatisticamente significativa ao nível de 95 \% entre as notas no componente de Conhecimento Específico dos diferentes tipos de Organização Acadêmica.
Sendo que a média dos concluintes das Faculdades ( 28,9 ) foi maior do que a dos concluintes de Centros Universitários ( 25,2 ) e de Universidades ( 26,9 ). 3.2 ANÁLISE DAS QUESTÕES OBJETIVAS 3.2.1 Componente de Formação Geral A Tabela 3.4 apresenta as estatísticas básicas relativas às oito questões objetivas do componente da prova que abrange a Formação Geral dos estudantes.
A média do Brasil foi 49,6.
A menor média foi encontrada na região Nordeste ( 49,0 ) e a maior na região Centro-Oeste ( 51,4 ).
As demais médias foram : 49,1 nas regiões Norte e Sudeste e 50,8 na região Sul.
O desvio padrão do Brasil foi 21,6, sendo o maior desvio padrão encontrado na região Sudeste ( 22,4 ) e o menor na região Sul ( 20,0 ).
Os demais desvios foram : 21,6 na região Nordeste, 21,4 na região Norte e 21,1 na região Centro-Oeste.
As mesmas medianas ( 50,0 ), notas máximas ( 100,0 ) e notas mínimas ( 0,0 ), nas questões objetivas do componente de Formação Geral, foram obtidas em todas as regiões, sem exceção.
A Tabela 3.5 apresenta os índices de facilidade e discriminação ( ponto bisserial ) das questões objetivas do componente de Formação Geral.
Quanto ao índice de facilidade, foram usadas as seguintes cores para diferenciar o nível de dificuldade da questão :  Azul para as questões classificadas com índice muito fácil ( índice \&gt ; =0,86 ), verde para as questões classificadas com índice fácil ( 0,61 a 0,85 ), amarelo para as questões classificadas com médio ( 0,41 a 0,60 ), vermelho para as questões classificadas com difícil ( 0,16 a 0,40 ) e roxo para as questões classificadas com muito difícil ( \&lt ; =0,15 ).
Já quanto ao índice de discriminação, foram usadas as seguintes cores para qualificar a questão :  As questões classificadas com índice fraco receberam a cor roxa ( índice \&lt ; =0,19 ), as classificadas com médio receberam a cor vermelho ( 0,20 a 0,29 ), as classificadas com bom receberam a cor verde ( 0,30 a 0,39 ) e as classificadas com muito bom ( \&gt ; =0,40 ) receberam a cor azul.
As questões objetivas do componente de Formação Geral, segundo o índice de facilidade, foram assim avaliadas : das oito questões, nenhuma teve o índice de facilidade classificado como muito fácil ou muito difícil.
Duas questões foram tidas como fáceis, por terem índice de acertos situado na faixa entre 0,61 e 0,85 ( de 61,0 \% a 85,0 \% de acertos ).
Quatro questões foram consideradas de dificuldade média, situando -se no intervalo entre 0,41 e 0,60 do índice de facilidade, ou seja, houve entre 41,0 \% e 60,0 \% de acertos.
Duas questões foram classificadas como difícil, situando -se no intervalo entre 0,16 e 0,40.
Como já comentado, para análise das questões objetivas relativas à Formação Geral segundo o poder de discriminação, utilizou -se, o índice de discriminação ponto bisserial.
Nesta análise as questões foram assim avaliadas : seis das oito questões apresentaram A Tabela 3.5 apresenta os índices de facilidade e discriminação ( ponto bisserial ) das questões objetivas do componente de Formação Geral.
Quanto ao índice de facilidade, foram usadas as seguintes cores para diferenciar o nível de dificuldade da questão :  Azul para as questões classificadas com índice muito fácil ( índice \&gt ; =0,86 ), verde para as questões classificadas com índice fácil ( 0,61 a 0,85 ), amarelo para as questões classificadas com médio ( 0,41 a 0,60 ), vermelho para as questões classificadas com difícil ( 0,16 a 0,40 ) e roxo para as questões classificadas com muito difícil ( \&lt ; =0,15 ).
Já quanto ao índice de discriminação, foram usadas as seguintes cores para qualificar a questão :  As questões classificadas com índice fraco receberam a cor roxa ( índice \&lt ; =0,19 ), as classificadas com médio receberam a cor vermelho ( 0,20 a 0,29 ), as classificadas com bom receberam a cor verde ( 0,30 a 0,39 ) e as classificadas com muito bom ( \&gt ; =0,40 ) receberam a cor azul.
As questões objetivas do componente de Formação Geral, segundo o índice de facilidade, foram assim avaliadas : das oito questões, nenhuma teve o índice de facilidade classificado como muito fácil ou muito difícil.
Duas questões foram tidas como fáceis, por terem índice de acertos situado na faixa entre 0,61 e 0,85 ( de 61,0 \% a 85,0 \% de acertos ).
Quatro questões foram consideradas de dificuldade média, situando -se no intervalo entre 0,41 e 0,60 do índice de facilidade, ou seja, houve entre 41,0 \% e 60,0 \% de acertos.
Duas questões foram classificadas como difícil, situando -se no intervalo entre 0,16 e 0,40.
Como já comentado, para análise das questões objetivas relativas à Formação Geral segundo o poder de discriminação, utilizou -se, o índice de discriminação ponto bisserial.
Nesta análise as questões foram assim avaliadas : seis das oito questões apresentaram índices acima de 0,40 e, assim, foram classificadas com índice muito bom para esse grupo de alunos ; uma questão teve índice de discriminação bom, entre 0,30 e 0,39.
Foi classificada com índice médio apenas uma das questões.
O índice de facilidade variou de 0,23 a 0,67, e o de discriminação, de 0,22 a 0,51.
As questões com índices de discriminação muito bom, de números 1, 2, 3, 4, 6 e 7, figuraram entre as mais fáceis desse conjunto : duas classificadas na categoria fácil ( questões 1 e 7 ) do índice de facilidade e quatro na categoria médio ( questões 2, 3, 4 e 6 ).
Em particular, as questões 2 e 7 foram as que apresentaram maior poder discriminatório, com índice 0,51, e foram também a terceira e segunda mais fáceis, com uma proporção de 0,56 e 0,62 de acertos, respectivamente.
A questão de número 5 apresentou índice de facilidade 0,23, classificando -a como difícil, ou seja, um quantitativo baixo de estudantes conseguiu resolvê -la, 23,0 \% dentro do universo de participantes.
Além disso, seu índice de discriminação foi médio ( 0,22 ).
Já a questão de número 8, também considerada difícil, obteve índice de discriminação 0,36, considerado bom.
O Gráfico 3.13, para exemplificar, analisa o comportamento da questão de número 7 de Formação Geral.
Trata -se da segunda questão mais fácil e uma das que obteve o maior índice de discriminação dessa parte da prova.
Neste gráfico, cada uma das cinco curvas representa o percentual de respostas em determinada alternativa da questão, em função da nota dos estudantes nesta parte da prova ( Formação Geral/Múltipla Escolha ), antes de possíveis eliminações pelo critério do ponto bisserial.
A curva em verde corresponde à alternativa B, a correta para esta questão.
Assim, observa -se que entre os estudantes com notas mais baixas, nessa parte do exame, a situação mais frequente foi a escolha de uma das alternativas incorretas : a alternativa C ( em preto ) ou A ( em azul ).
Na medida em que a nota aumenta, indicando desempenho melhor nesta parte da prova, aumenta concomitantemente a proporção de estudantes que selecionaram a alternativa correta B, atingindo 100 \% para os alunos com 8 acertos.
Essa análise permite verificar como a questão discriminou os grupos de desempenho, justificando o alto índice obtido na questão. situação mais frequente foi a escolha de uma das alternativas incorretas : a alternativa C ( em preto ) ou A ( em azul ).
Na medida em que a nota aumenta, indicando desempenho melhor nesta parte da prova, aumenta concomitantemente a proporção de estudantes que selecionaram a alternativa correta B, atingindo 100 \% para os alunos com 8 acertos.
Essa análise permite verificar como a questão discriminou os grupos de desempenho, justificando o alto índice obtido na questão.
Os gráficos relativos às demais questões de Formação Geral constam do Anexo I. 3.2.2 Componente de Conhecimento Específico A Tabela 3.6 apresenta as estatísticas básicas em relação às questões objetivas do componente de Conhecimento Específico da prova, por Grande Região.
A média do Brasil deste componente foi de 29,4.
A menor média foi observada na região Nordeste ( 27,2 ) e a maior na região Sudeste ( 30,3 ).
O desvio padrão de todo o Brasil foi 13,8, sendo o menor desvio padrão encontrado na região Norte ( 10,7 ) e o maior na região Centro-Oeste ( 15,1 ).
A mediana de todo o Brasil foi 29,2.
Na maioria das regiões do Brasil, exceto as regiões Norte e Nordeste ( 25,0 ), a mediana também foi 29,2.
A nota máxima da prova ( 95,8 ) foi obtida, nas questões objetivas do Componente Específico, por pelo menos um aluno da região Sudeste.
A nota máxima nas demais regiões foi : 91,7 na região Centro-Oeste, 83,3 na região Sul, 75,0 na região Nordeste e 70,8 na região Norte.
A nota mínima foi zero em todas as regiões, sem exceção.
A Tabela 3.7 apresenta os índices de facilidade e discriminação ( ponto bisserial ) das questões objetivas do componente de Conhecimento Específico.
Para facilitar a diferenciação das questões usaram -se as mesmas cores da Tabela 3.5 para as diferentes classificações dos índices de facilidade e de discriminação.
Dentre as questões objetivas da parte da prova relativa ao componente de Conhecimento Específico, nenhuma foi anulada pela Comissão.
Desse modo, a classificação quanto ao índice de facilidade foi estabelecida em todas as 27 questões.
A partir dos índices obtidos, pode -se concluir que a maioria das questões objetivas da prova foi considerada pelo menos difícil : das 27 questões, 18 foram classificadas como difíceis e quatro como muito difíceis.
Não houve questão classificada como muito fácil ou fácil, e cinco questões consideradas médias, entre 0,41 e 0,60.
Já quanto aos índices de discriminação das questões objetivas do componente de Conhecimento Específico da prova, tem -se como resultado a seguinte classificação : dez das 27 questões foram consideradas como boas, enquanto duas delas tiveram índice de discriminação muito bom.
Assim, para 12 das 27 questões, os índices de discriminação foram bons ou muito bons.
Dentre as demais, 12 delas foram classificadas como médias e outras três como fracas, sendo 15, por conseguinte, a quantidade de questões nos dois patamares mais baixos de discriminação.
Constata -se, assim, que a prova – no que se refere -se ao componente de Conhecimento Específico – possuía capacidade razoável de discriminar entre aqueles que dominam ou não o conteúdo.
Dentre as questões que alcançaram os maiores índices de discriminação, as de números 12 e 16 classificadas com índice muito bom, situando -se no intervalo de 0,40 a 0,45 do índice, ambas foram classificadas na categoria média, quanto ao índice de facilidade.
A questão de número 28 foi a mais difícil dentre as 27 questões específicas, com baixo índice de facilidade, apenas 9,0 \% de acertos.
Essa questão apresentou poder discriminatório igualmente baixo, 0,09, o que comprova ter sido esta a mais difícil para os estudantes.
Destaca -se, também, a questão 35, com índice de facilidade 0,14, o que, em termos percentuais, corresponde a 14,0 \% de estudantes que responderam acertadamente, obtendo, ainda, 0,11 de índice de discriminação.
Tais questões foram, portanto, duas das mais difíceis da prova.
As questões 28 e 35 foram eliminadas do cômputo da nota final pelo critério do ponto bisserial.
Além destas duas, a questão 17, igualmente com índice fraco de discriminação, também não foi computada.
A título de exemplo das análises do comportamento das questões objetivas, o Gráfico 3.14 analisa a questão 16 do componente de Conhecimento Específico.
Esta foi a segunda questão mais fácil da prova, apresentando índice de facilidade 0,50, ou seja, 50,0 \% dos estudantes assinalaram acertadamente a opção D, correspondente ao gabarito.
Seu índice de discriminação foi igual a 0,42, classificado como muito bom.
Neste gráfico, cada uma das cinco curvas representa o percentual de respostas em determinada alternativa da questão 16, em função da nota dos estudantes nesta parte da prova, antes de possíveis eliminações de questões pelo critério do ponto bisserial.
A alternativa correta D, representada no gráfico pela curva em roxo, foi escolhida em maiores proporções pelos alunos com desempenho melhor nesta parte da prova.
Já as alternativas incorretas, também denominadas distratores, foram selecionadas principalmente por aqueles com notas mais baixas.
Aqueles com nota zero, na sua quase totalidade deixaram esta questão em branco ou marcaram mais de uma alternativa, comportamento considerado inválido.
A proporção de alunos que selecionou a resposta correta D aumenta gradativamente, chegando a atingir 100 \% para 19 acertos e a partir de 21 acertos, enquanto a proporção dos que escolheram alternativas incorretas decai, a partir da primeira nota não nula, como função da nota nesta parte da prova.
Os gráficos relativos às demais questões do componente de Conhecimento Específico constam do Anexo I. 3.3 ANÁLISE DAS QUESTÕES DISCURSIVAS 3.3.1 Componente de Formação Geral As análises dos resultados de desempenho dos estudantes de Ciências Econômicas nas duas questões discursivas relativas à Formação Geral encontram -se na Tabela 3.8 e no Gráfico 3.15.
Na tabela 3.8 observa -se que as notas médias foram mais baixas nesse conjunto de questões do que no das objetivas.
Os estudantes de todo o Brasil obtiveram, em Formação Geral, média 49,6 nas questões objetivas e 33,9 nas questões discursivas.
A mediana também confirma o pior desempenho dos alunos de todo o Brasil nas questões discursivas do componente de Formação Geral.
Enquanto essa estatística foi de 40,0 para questões discursivas, para as questões objetivas foi de 50,0.
Pode -se, também, notar um aumento do desvio padrão de 21,6, nas questões objetivas do componente de Formação Geral dos alunos de todo o Brasil, para 25,5 nas questões discursivas do mesmo componente.
Como já comentado, a mediana de todo o Brasil, neste componente, foi 40,0, sendo a maior mediana encontrada na região Norte ( 45,0 ) e a menor nas regiões Nordeste e Sudeste ( 37,5 ).
A nota máxima ( 95,0 ) foi obtida, nas questões discursivas do componente de Formação Geral, por pelo menos um aluno na região Sudeste.
A nota máxima nas demais regiões foi : 90,0 na região Centro-Oeste, 87,5 na região Sul, 85,0 na região Nordeste e 82,5 na região Norte.
A nota mínima foi zero em todas as regiões, sem exceção.
O Gráfico 3.15 representa a distribuição das notas nas questões discursivas no componente de Formação Geral.
A moda desta distribuição ocorre no primeiro intervalo, [ 0 ; 10 ], devido à grande quantidade de notas zero e à alta frequência de alunos que deixaram este tipo de questão em branco.
Desconsiderando esse intervalo, a distribuição possui outro intervalo modal, ( 50 ; 60 ].
Destaca -se, também, o intervalo ( 40 ; 50 ] que possui frequência muito próxima do intervalo ( 50 ; 60 ].
A distribuição possui assimetria à esquerda, com coeficiente de assimetria -0,13.
Em todas as regiões o coeficiente de assimetria é negativo ( entre -0,42 para a região Norte e -0,07 para a região Sudeste ).
Na sequência, os resultados verificados para cada uma das questões discursivas de Formação Geral serão apresentados, estabelecendo -se relações com os conteúdos abordados em cada uma delas.
Os comentários da Banca de docentes corretores a respeito do observado na correção das respostas dos estudantes, suas impressões e conclusões serão apresentados junto à análise de cada questão.
Cumpre esclarecer que, tendo em vista que as questões discursivas de Formação Geral são padronizadas, ou seja, constam de todas as provas, os comentários da Banca são os mesmos para todas as carreiras acadêmicas, sendo direcionados a todos os estudantes que participaram do ENADE/2012.
A seguir, serão analisados os desempenhos da Área de Ciências Econômicas nas duas questões discursivas de Formação Geral do ENADE/2012, comparando os resultados obtidos com comentários para cada questão. 3.3.1.1 Análise da Questão Discursiva 1 do Componente de Formação Geral Os dados de Ciências Econômicas, obtidos a partir das respostas à questão 1, encontram -se na Tabela 3.9 e no Gráfico 3.16.
Nessa questão – de pior desempenho dentre as duas de Formação Geral – os alunos de todo Brasil tiveram média 32,7.
A maior média para a questão 1 foi obtida na região Norte ( 37,8 ), e a menor, na região Nordeste ( 31,6 ).
Quanto à variabilidade das notas, o desvio padrão de todo o Brasil foi 27,7.
O menor desvio padrão foi obtido na região Sul ( 26,8 ) e o maior nas regiões Nordeste e Sudeste ( 28,0 ).
A mediana para o Brasil como um todo foi 30,0.
A maior mediana foi obtida nas regiões Norte e Centro-Oeste ( 40,0 ), enquanto a menor mediana ocorreu nas regiões Nordeste e Sudeste ( 30,0 ).
A nota máxima ( 95,0 ) da questão discursiva 1 foi a mesma para todas as regiões do Brasil, sem exceção.
Além disso, a mesma nota mínima ( 0,0 ) foi obtida em todas as regiões do país, também sem exceção.
O Gráfico 3.16 mostra a distribuição das notas na questão discursiva 1 do componente de Formação Geral.
Observa -se que a maior frequência corresponde aos alunos que deixaram a questão em branco.
Considerando apenas os estudantes que tentaram resolver a questão, a moda ocorre no intervalo ( 40 ; 50 ].
Também observa -se grande concentração de notas no intervalo ( 30 ; 40 ].
O coeficiente de assimetria da distribuição das notas é positivo ( 0,18 ), indicando uma assimetria à direita.
Em todas as regiões o coeficiente de assimetria é positivo ( entre 0,04 para a região Centro-Oeste e 0,22 para as regiões Nordeste e Sudeste ), exceto na região Norte onde o coeficiente foi negativo ( -0,07 ). 3.3.1.2 Comentários sobre a correção das respostas à Questão Discursiva 1 A questão aborda um tema bastante polêmico o qual vem sendo amplamente discutido, a conciliação entre o desenvolvimento econômico e a preservação dos recursos naturais.
Busca -se observar a relação entre o estímulo governamental à compra de veículos novos e os impactos ambientais decorrentes desta política pública.
Sua compreensão implica no conhecimento de uma temática atual e na reflexão critica a respeito das decisões governamentais em sua correlação com o cenário econômico internacional e os problemas e soluções de ordem ambiental.
Embora o comando da questão solicitasse explicitamente que a resposta fosse fornecida em um texto único ( “ redija um texto dissertativo sobre... ” ), foi significativamente alta foi a proporção de respondentes que atendeu a cada item ( a, b e c ) do enunciado de maneira isolada, como partes estanques da mesma resposta.
No que se segue encontra -se uma avaliação da correção de cada um dos três itens. a - “ conceito de desenvolvimento sustentável ” A variedade de respostas a esse item confirmou a importância do tema para avaliar a formação geral dos estudantes e cursos, já que ( i ) é atual ; ( ii ) foi abordado nos meios de comunicação ( especialmente no ano de 2012, em que o Brasil sediou a “ Rio+20 ” ) ; ( iii ) é relevante não apenas no Brasil como em todos os principais fóruns internacionais ; e ( iv ) é comum a diversas áreas de conhecimento.
Outros dois aspectos positivos merecem destaque : ( i ) a grande incidência de respostas que apontaram a relação entre gerações como aspecto essencial da ideia de sustentabilidade socioambiental – no mais, as respostas apresentaram pelo menos a tentativa de indicar uma definição de desenvolvimento sustentável ou a abordagem positiva dessa mesma noção, evidenciado compreensão do enunciado ; e ( ii ) o exercício crítico que a questão propunha, no sentido de exigir dos respondentes a compreensão das contradições de políticas públicas e a apresentação de propostas de soluções para conciliar o incentivo do crescimento econômico e a promoção da sustentabilidade socioambiental, a partir de uma perspectiva transformadora da realidade.
Apesar dos aspectos positivos apontados no parágrafo anterior, a maioria dos estudantes não alcançou o máximo da pontuação.
De forma geral, trataram o desenvolvimento sustentável como sinônimo de não agressão ambiental.
Alguns confundiram desenvolvimento sustentável com desenvolvimento econômico, e outros faziam referência apenas à preocupação com a conservação dos recursos mencionando a preocupação com gerações futuras.
Deve -se, ainda, destacar o significativo número de alunos que apresentou uma definição aparentemente copiada do enunciado de outra questão da prova, como por exemplo, a questão objetiva 31 da prova de Administração. b - “ conflito entre o estímulo à compra de veículos automotores e a promoção da sustentabilidade ” Esse foi o item identificado como aquele no qual as respostas dos estudantes avaliados mais se aproximaram da pontuação máxima prevista no padrão de resposta.
Ou seja, a ampla maioria dos respondentes foi bem-sucedida em apontar a relação entre o aumento da quantidade de veículos automotores em circulação e os danos socioambientais que isso acarreta no cenário urbano, devido ao aumento da emissão de gases poluentes.
No item b, as respostas corretas se concentraram em críticas ao incentivo oferecido pelo governo, referências ao problema dos congestionamentos e ao aumento da poluição.
Contudo, a maioria das respostas não explicitava a correlação entre o aumento da circulação de veículos automotores e o agravamento dos problemas ambientais.
Um número expressivo de alunos apresentou uma resposta em que era dado destaque à política de redução do IPI como um mecanismo para renovação da frota de automóveis ou para a aquisição do transporte próprio e da consequente redução da dependência ao transporte coletivo.
Também foram notados alguns tipos de padrões recorrentes nas tentativas de resposta a este item.
Dois desses padrões se destacaram : ( i ) as declarações genéricas de condenação do “ consumismo ” de uma “ nova classe média ” ou da busca pelo “ status ” que a posse de um automóvel particular proporcionaria, como se fosse possível atribuir a essa ordem de coisas a causa principal do aumento dos congestionamentos nas grandes cidades ; e ( ii ) a repetição da avaliação, presente no enunciado da questão, que aponta de forma superficial a contradição entre o estímulo à demanda por automóveis particulares e a promoção da sustentabilidade. c - “ ações de fomento ao transporte urbano sustentável no Brasil ” Este foi o item que ofereceu maior facilidade de aplicação do padrão de resposta, mas foram poucos os casos de pontuação integral.
Dentre as soluções mais destacadas pelos alunos apareceram : ( i ) o incentivo ao uso de energias alternativas ; e ( ii ) o estímulo ao transporte coletivo ( metrô, trem e ônibus ), além do uso de bicicletas.
Importante destacar o fato de que muitos apresentavam respostas com referências à necessidade de incentivo ao transporte urbano, nas quais este termo era utilizado como sinônimo de transporte coletivo.
A avaliação da questão 1 trouxe à tona problemas estruturais da educação básica, que se refletem no ensino superior.
Foi muito grande o número de provas em que o aluno demonstrou não compreender o enunciado da questão, apresentando respostas que fugiam ao tema proposto, mesmo quando bem redigidas.
Uma resposta recorrente fazia referência às consequências da ampla oferta de crédito à população e aos problemas relacionados à inadimplência, situação que, embora, verdadeira, não estava relacionada à questão proposta.
Houve, ainda, os casos em que as respostas apresentadas não só se afastaram do tema como não reuniam qualquer tipo de argumentação coerente.
E, por fim, destaca -se a existência de respostas incoerentes e desconexas.
Outro aspecto geral diz respeito ao grande número de provas em branco ou em que o aluno apresentou a resposta a outro item no espaço da questão 1.
Considerando -se o objetivo geral deste exame, avaliar a qualidade dos cursos de graduação oferecidos no país e contribuir para a melhoria da qualidade do ensino superior, destaca -se a imensa responsabilidade do processo que envolve a elaboração deste diagnóstico.
Nesse sentido, é urgente alertar para a ocorrência de um número significativo de respostas que se afastam de um padrão mediano de qualidade esperado de um aluno concluinte em um curso de graduação, especialmente em se considerando o fato de que a questão proposta faz referência a tema intensamente explorado pela mídia. 3.3.1.3 Análise da Questão Discursiva 2 do Componente de Formação Geral A Tabela 3.10 mostra que o desempenho dos estudantes na questão 2 ( média 35,2 ) foi superior ao obtido na questão de número 1 ( média 32,7 ).
A região Centro-Oeste foi aquela onde a média, nessa questão, foi maior ( 40,0 ), e a de menor média foi a região Sudeste ( 34,1 ).
Quanto à variabilidade das notas, o desvio padrão de todo o Brasil foi 29,1, superior ao obtido na questão de número 1 ( 27,7 ).
O maior desvio nessa questão foi obtido na região Sudeste ( 29,5 ), enquanto o menor foi obtido na região Sul ( 28,1 ).
A mediana para o Brasil como um todo foi 45,0.
A menor mediana ocorreu nas regiões Nordeste e Sudeste ( 40,0 ) e a maior ocorreu na região Centro-Oeste ( 50,0 ).
A nota máxima ( 95,0 ) foi obtida nas regiões Sudeste e Sul, já a menor nota máxima foi obtida nas regiões Norte e Centro-oeste.
A mesma nota mínima ( 0,0 ) foi obtida em todas as regiões do país, sem exceção.
O Gráfico 3.17 mostra a distribuição das notas na questão discursiva 2 do componente de Formação Geral.
Neste gráfico, desconsiderando -se a grande quantidade de alunos que deixaram a questão 2 em branco ( em torno de 34 \%, correspondendo à moda ), observa -se a moda no intervalo ( 50 ; 60 ].
O coeficiente de assimetria das notas nesta questão é negativo ( -0,08 ), indicando que a distribuição é assimétrica à esquerda.
O mesmo ocorre para todas as regiões do Brasil, cujos coeficientes variam entre -0,34 para a região Centro-Oeste e -0,01 para a região Sudeste. 3.3.1.4 Comentários sobre a correção das respostas à Questão Discursiva 2 O conteúdo da questão era bastante adequado às Diretrizes do INEP para o ENADE, além de atual e de abrangente.
Porém, exigia -se a articulação de muitas variáveis.
Buscando aferir, a partir da apresentação de uma charge e de um texto introdutório, as competências e habilidades relacionadas à leitura, à interpretação, à análise crítica, à capacidade argumentativa, à elaboração de sínteses e de avaliação prospectiva sobre tema da atualidade, a questão tratava da violência e suas repercussões no contexto atual.
O comando da questão solicitava, a partir da definição de violência da OMS e de duas imagens, que o aluno elaborasse um texto, de caráter dissertativo, sobre a “ violência na atualidade ”.
Embora tal solicitação fosse explícita ( “ redija um texto dissertativo sobre... ” ), como na questão 1, foi alta a proporção de alunos que respondeu a cada item do enunciado isoladamente, como partes independentes. capacidade argumentativa, à elaboração de sínteses e de avaliação prospectiva sobre tema da atualidade, a questão tratava da violência e suas repercussões no contexto atual.
O comando da questão solicitava, a partir da definição de violência da OMS e de duas imagens, que o aluno elaborasse um texto, de caráter dissertativo, sobre a “ violência na atualidade ”.
Embora tal solicitação fosse explícita ( “ redija um texto dissertativo sobre... ” ), como na questão 1, foi alta a proporção de alunos que respondeu a cada item do enunciado isoladamente, como partes independentes.
Verificaram -se poucas respostas em branco, como também foram poucas as respostas contendo protestos contra o ENADE, muito frequentes em anos anteriores, o que significa um avanço para a afirmação do ENADE como um exame aceito socialmente. a -"tecnologia e violência"Tanto as charges quanto o texto apresentado no enunciado induziram os alunos a focarem suas respostas nos meios de comunicação.
Além disso, algumas respostas estabeleciam a tecnologia, em si, como causa da violência tanto nas escolas ( em particular ), quanto na sociedade ( no geral ).
Muitos respondentes tiveram dificuldade para fazer a relação entre tecnologia, violência e a escola, de modo adequado, centrando -se na violência nas escolas.
Além disso, a grande maioria das respostas girava em torno da exposição da violência através da TV, internet e jogos digitais, sendo que predominaram as abordagens com ênfase em programas e filmes veiculados pela televisão.
Poucos foram os estudantes que estabeleceram relação entre o contexto socioeconômico e a violência. b -"causas e consequências da violência na escola"Embora a solicitação desse item fosse objetiva e atendida por grande parte dos alunos, inúmeras vezes, observou -se confusão entre o significado de causa e de consequência.
Em outras inúmeras vezes, as respostas apontavam apenas para uma das classes, causas ou consequências.
Foram frequentes as respostas que atribuíam o aumento da violência à ineficiência do poder público, apontando a falta de investimentos, o baixo salário dos professores, a ausência de profissionais qualificados como causas.
Foram frequentes também os que apontavam a violência doméstica, o bullying, e a transferência do imaginário dos filmes e jogos para a realidade, como causas da violência. c -"proposta de solução para o problema da violência na escola"Sem dúvida, há uma grande diversidade de abordagens possíveis para apresentação de uma proposta de solução para o problema da violência, o que o padrão de resposta buscou contemplar.
No entanto, houve grande número de respostas baseadas em opiniões pessoais dos estudantes e vinculadas ao senso comum.
As respostas giraram em torno da necessidade de mais investimentos na educação por parte do poder público e seu rebatimento no maior aparelhamento das escolas, salário dos professores etc.
Além disso, foram muitas as respostas que trataram da necessidade de maior conscientização dos pais, alunos e professores com a realização de um trabalho conjunto que envolvesse todos os segmentos da sociedade.
De um modo geral, os estudantes demonstraram dificuldades em relação ao nível de abrangência exigido, falta de conteúdo, inabilidade na elaboração de texto escrito.
Da grande quantidade de provas corrigidas, o que ficou evidente, foi a dificuldade de estabelecer as conexões requeridas na questão 2.
A violência nas escolas foi bastante associada ao reflexo do meio onde a criança e o adolescente mantém os seus vínculos afetivos.
Surgiram muitas respostas revelando críticas à intolerância no núcleo familiar, à jornada de trabalho dos pais, ao distanciamento da mãe nas atividades domésticas etc.
Em boa parte das respostas, os estudantes lançaram a causa da violência escolar exclusivamente na ausência de participação familiar na vida dos filhos, às vezes criando uma cisão com a proposta da questão, já que abandonavam totalmente a temática da tecnologia em suas reflexões.
Além da ausência dos pais no cotidiano da escola, divórcios e atos de violência contra crianças foram considerados fatores que delegariam para a escola novas funções.
Também foram mencionadas a má formação de professores e a insatisfação com o ambiente escolar.
Em inúmeros casos, ficou evidente o anseio de se fazer da escola um lugar de resgate tanto da criança e do adolescente quanto da família, compreendendo -se a escola em sentido social e assistencial.
Segundo este amplo grupo de concluintes, a escola deveria oferecer -se como centro de cultura e lazer, como espaço para o acompanhamento psicológico e serviço social às comunidades diretamente envolvidas.
Ressalte -se o constante aparecimento das sugestões de horário integral, melhores salários para professores e iniciativas para melhoria na formação de professores.
As propostas de solução para acabar com a violência na escola se detiveram em aspectos muito gerais e poucos alunos detalharam uma proposta de intervenção, apresentando ações concretas.
Dentre as propostas de solução surgiram : educação familiar, ensino religioso, maior segurança nas escolas, com utilização de aparelhos para detecção de armas, e o retorno do papel do professor como um profissional respeitado pelos poderes municipais e estaduais.
Um percentual significativo de provas apresentou resposta simplista, com ênfase em determinados pontos como : religião ( sem demonstração do conceito de laicidade ), políticas sociais ( sem demonstração do domínio de conceitos que envolvem Estado, Governo e Politicas Públicas ) etc.
Foi grande a incidência de respostas que se referiam às leis, à vigilância e maior controle dos pais e professores ; no entanto, não demonstram entender a noção de legalidade para os diferentes segmentos sociais e como os mesmos se encontram submetidos a diferentes formas de poder.
Alguns alunos, ainda que raros, trataram a questão com a profundidade esperada de concluintes de cursos universitários.
Por outro lado, foram vários os textos que apresentaram uma argumentação simplória e desconexa, com forte visão paternalista sobre o papel do governo e do professor, como se tais instâncias pudessem intervir na educação familiar.
Enquanto alguns conseguiram desenvolver o tema proposto com elegância e desenvoltura, encadeando o tema de forma harmônica, com excelente nível de linguagem escrita e correção gramatical, outros sequer conseguiram conectar as ideias e construir um texto razoável, utilizando, inclusive, um linguajar mais propício à linguagem falada do que à escrita.
Cabe também ressaltar a quantidade de respostas carregadas de conteúdo com juízo de valor discriminatório ou extremamente conservador, do tipo : “ mulher não deve trabalhar fora ” ; “ amar e temer a Deus é a solução ” ; obrigatoriedade do ensino religioso e de uma doutrina religiosa ; imposição de uma educação do tipo militar nas escolas ; censura nos meios de comunicação etc.
Outras soluções apontadas para resolver o problema da violência nas escolas estavam associadas à, por exemplo : redução da maioridade penal, defesa da pena de morte, milícias armadas na porta das escolas, punições como o retorno da palmatória.
O Estatuto da Criança e do Adolescente foi uma peça legal criticada, e recebeu, em diversas respostas, o título de responsável pelo sentimento de impunidade de adolescentes infratores.
Contudo, algumas dessas proposições, desde que bem contextualizadas e formuladas com boa argumentação, foram pontuadas de forma a respeitar o direito democrático de pluralidade opinativa.
É imperioso ressaltar que a falta de inteligibilidade segue representando um aspecto merecedor de providências na formação universitária, uma vez que exprime uma habilidade importante para qualquer profissional de nível superior.
Além de terem sido poucas respostas boas e claras, foi possível identificar a deficiência de argumentação, além dos inúmeros problemas gramaticais ( ortografia, concordância, pontuação, regência, etc. ).
Salienta -se, nesse sentido, o uso indiscriminado de letras maiúsculas e minúsculas numa mesma frase, sem obedecer ao critério formal da língua portuguesa. 3.3.2 Componente de Conhecimento Específico Na parte da prova relativa às questões discursivas no componente de Conhecimento Específico ( Tabela 3.11 ), observa -se que a média foi mais baixa do que para as questões discursivas do componente de Formação Geral.
Enquanto no componente de Formação Geral a média para estudantes de Ciências Econômicas de todo o Brasil foi 33,9, na parte de Conhecimento Específico a média foi 13,6.
A maior média deste componente foi obtida pelos estudantes da região Sudeste ( 14,9 ), e a menor, pelos da região Nordeste ( 10,8 ).
Quanto à variabilidade das notas, o desvio padrão de todo o Brasil foi 17,9.
Mantendo o mesmo comportamento da média, o maior desvio padrão foi encontrado na região Sudeste ( 19,2 ), e o menor, na região Nordeste ( 15,2 ).
A maior nota máxima, nas questões discursivas do componente de Conhecimento Específico, foi obtida na região Sudeste ( 96,7 ) enquanto a menor nota máxima foi encontrada na região Norte ( 80,0 ).
Além disso, a mesma nota mínima ( 0,0 ) foi obtida por alunos de todas as regiões do Brasil sem exceção.
A mediana do Brasil como um todo foi 5,0, também encontrada nas regiões Sudeste e Centro-Oeste.
As demais medianas foram : 6,7 na região Sul, 3,3 na região Norte e 0,0 na região Nordeste.
O Gráfico 3.18 representa a distribuição das notas nas questões discursivas no componente de Conhecimento Específico.
Essa distribuição tem moda no intervalo de [ 0 ; 10 ].
Nota -se que a partir desse intervalo a distribuição é decrescente até o último intervalo.
O coeficiente de assimetria é positivo e maior do que 1 ( 1,40 ), indicando uma grande assimetria à direita.
Este coeficiente também é positivo em todas as regiões do Brasil, variando entre 1,28 para a região Sudeste e 1,61 para a região Norte. 3.3.2.1 Análise da Questão Discursiva 3 do Componente de Conhecimento Específico Na questão 3, cujos resultados aferidos encontram -se descritos na Tabela 3.12, a média dos estudantes de todo o Brasil foi 10,0.
A menor média nessa questão foi obtida pelos alunos da região Nordeste ( 7,3 ), enquanto a maior média foi obtida na região Sudeste ( 11,4 ).
Quanto à variabilidade das notas, o desvio padrão de todo o Brasil foi 20,7.
O maior desvio padrão foi obtido na região Sudeste ( 22,5 ), enquanto o menor foi obtido na região Nordeste ( 16,9 ).
A nota máxima, 100,0 pontos, foi alcançada por pelo menos um aluno em todas as regiões, sem exceção.
As mesmas medianas ( 0,0 ) e notas mínimas ( 0,0 ) foram obtidas em todas as regiões do Brasil, também sem exceção.
Ou seja, pelo menos metade dos alunos de cada uma das regiões tirou nota zero nesta questão.
O Gráfico 3.19 mostra a distribuição das notas na questão discursiva 3, do componente de Conhecimento Específico.
Essa distribuição é unimodal, com moda nas questões em branco ( também com nota zero ).
Considerando -se em separado a distribuição das notas dos alunos que resolveram a questão 3, observamos que o gráfico é fortemente assimétrico à direita, com alto coeficiente de assimetria positivo 2,33.
Em todas as regiões do Brasil o coeficiente de assimetria foi positivo ( entre 2,16 para a região Sudeste e 2,75 para a região Norte ). 3.3.2.2 Comentários sobre as respostas à Questão Discursiva 3 A questão se insere no conteúdo programático na parte de Teoria Econômica do curso de Ciências Econômicas, mais especificamente na parte de Teoria Microeconômica, sendo que a letra ( b ) do enunciado envolve também conhecimentos ligados à disciplina de Economia Internacional.
O enunciado é claro e as perguntas objetivas, envolvendo conhecimento de duas disciplinas do curso de graduação que são fundamentais para a boa formação do Economista.
O processo de correção da questão 3 teve como uma das principais características a quase ausência de elementos gráficos para apoiar as respostas.
Além disso, foi difícil identificar nas respostas dos discentes uma articulação ( exigida no item ( b ) da questão ) entre os campos de microeconomia e economia internacional.
Para o item ( a ), a correção considerou aproximações dos termos usados no padrão de \textbf{reposta} aceitando expressões que davam entendimento de que o discente entendia questões relacionadas à Fronteira de Possibilidade de Produção.
Já no item ( b ), não encontradas muitas variações entre as respostas oferecidas e o padrão de resposta.
A questão possui o mérito de exigir do discente uma boa formação na Área de microeconomia e sua interação com outro campo de conhecimento no curso de economia. 3.3.2.3 Análise da Questão Discursiva 4 do Componente de Conhecimento Específico A Tabela 3.13 contém as informações relativas à questão 4 do conjunto de questões do componente de Conhecimento Específico.
O desempenho dos estudantes de todo o Brasil nesta questão foi superior ao desempenho na questão de número 3.
A média geral do Brasil foi 19,9, sendo a menor média registrada na região Nordeste ( 17,8 ) e a maior na região Norte ( 21,0 ).
A nota máxima ( 100,0 ) foi atingida em todas as regiões, sem exceção.
A mesma mediana ( 0,0 ) e a mesma nota mínima ( 0,0 ) foram também obtidas em todas as regiões do Brasil, sem exceção.
Ou seja, pelo menos metade dos alunos em cada uma das regiões tirou nota zero nesta questão.
O Gráfico 3.20 representa a distribuição de notas na questão discursiva 4, no componente de Conhecimento Específico.
Como nas outras questões discursivas, o número de estudantes que deixou a questão em branco foi muito grande, correspondendo à moda da distribuição ( acima de 45 \% ).
Dentre as notas dos que resolveram a questão 4, o intervalo modal foi ( 10 ; 20 ].
A assimetria da distribuição confirma o valor baixo da média, que foi 19,9, já que apresenta concentração em notas inferiores ao intervalo ( 10 ; 20 ] e uma cauda longa do lado direito deste intervalo.
Tanto na distribuição como um todo quanto na análise dos resultados dos alunos por região, os coeficientes de assimetria são sempre positivos e com valores maiores do que 1.
No Brasil como um todo o coeficiente de assimetria foi 1,09 e nas regiões variou entre 1,03 para a região Norte e 1,25 para a região Nordeste. 3.3.2.4 Comentários sobre as respostas à Questão Discursiva 4 O tema da questão faz parte do bloco de formação econômica do Brasil e economia brasileira, sendo bastante atual.
A questão cobrava do discente uma comparação entre a crise de 1929 e a recente crise de 2008.
Além disso, tinha o mérito de exigir do discente, também, conhecimentos de ferramentas de macroeconomia, já que o enfrentamento das crises passou pelo uso de instrumentos desse campo do conhecimento.
Durante a correção foi identificado que os discentes possuem, em sua grande maioria, dificuldades em estabelecer uma conexão entre o processo histórico e as ferramentas de intervenção do economista na tentativa de solucionar problemas gerados por crises circunstanciais.
Vale dizer que, em alguma medida, a questão exigia conhecimento básico de ferramentas de macroeconomia do bloco de teoria econômica do curso de economia.
Muitas respostas foram dominadas pela ideia de que, no caso da economia brasileira, nos dois períodos, a superação veio através de estímulos ao consumo doméstico, com políticas expansionistas, sem mencionar a importante diferença quando se fala da crise de 29 e o importante efeito desta sobre o movimento de deslocamento produtivo observado na economia brasileira. 3.3.2.5 Análise da Questão Discursiva 5 do Componente de Conhecimento Específico A Tabela 3.14 contém as informações relativas à questão 5 do conjunto do componente de Conhecimento Específico.
O desempenho dos estudantes nessa questão foi inferior ao da questão 4 e superior ao da questão 3.
A nota média dos estudantes de todo o Brasil foi 10,8.
A maior média foi registrada na região Sudeste ( 12,9 ), enquanto a menor média foi registrada na região Norte ( 6,9 ).
Quanto à variabilidade das notas, o desvio padrão dos alunos do Brasil, como um todo, foi 21,8.
Enquanto o maior desvio foi encontrado na região Sudeste ( 23,8 ), o menor foi encontrado na região Norte ( 17,8 ).
A mesma nota máxima ( 100,0 ) foi alcançada em todas as regiões, sem exceção.
A mediana e a nota mínima foram zero para todas as regiões do Brasil, sem exceção.
Ou seja, pelo menos metade dos alunos tirou nota zero nesta questão.
O Gráfico 3.21 representa a distribuição das notas da questão discursiva 5 do componente de Conhecimento Específico.
Destaca -se, como nas demais questões discursivas, o grande número de estudantes que deixaram esta questão em branco, correspondendo à moda.
O coeficiente de assimetria para todos os alunos é 2,20, indicando uma grande assimetria à direita, e para todas as regiões se mantém positivo, variando entre 1,91 para a região Sudeste e 3,02 para a região Norte. 3.3.2.6 Comentários sobre as respostas à Questão Discursiva 5 O enunciado da questão estava claro e objetivo, deixando pouco espaço para interpretações diferentes do que seria exigido como padrão de resposta.
Comparando os dois itens da questão, verifica -se um grau de dificuldade maior no item ( b ), porém nada que estivesse em desacordo com o padrão de conhecimento exigido pelas Diretrizes Curriculares do curso de Ciências Econômicas.
Assim como as demais, a questão tinha o mérito de exigir do discente uma capacidade de interagir com campos de conhecimento distintos do curso de economia.
No caso dessa questão a interação entre economia brasileira e macroeconomia.
No item ( a ) observou -se dificuldade em estabelecer uma articulação entre a política cambial do período e seu principal objetivo ( estabilização monetária ).
Em relação ao item ( b ) observou -se um enorme desconhecimento sobre a estrutura do balanço de pagamentos de um país e a articulação entre déficit em conta corrente e sua contrapartida na conta de capitais.
A questão, além de bem elaborada, era bastante contemporânea no que diz respeito à formação do economista brasileiro, pois é possível verificar questões semelhantes nos principais concursos públicos para o cargo de economista no país. 3.3.3 Considerações Finais As médias das questões deixam claro que em termos nacionais o desempenho geral dos discentes ficou abaixo do que se espera de um profissional de economia bem formado.
Um dado da correção que ficou claro foi a quase completa ausência de habilidade dos discentes em lidar com representações gráficas ( exigida particularmente na questão 3 ), além de um enorme desconhecimento sobre a estrutura de balanço de pagamento de uma economia.
Outro elemento que ficou claro foi o desconhecimento teórico do modelo de estabilização que serviu de base ao Plano Real, o que em alguma medida mostra desconhecimento também de fundamentos de macroeconomia.
Em relação à questão de maior \textbf{peso} no bloco de História econômica e formação econômica do Brasil ( questão 4 ), as respostas demonstraram desconhecimento tanto do processo histórico quanto das decisões de política econômica implementadas.

% matched lemmas: peso, repor
\end{textsample}
